% Caso de Uso: Eliminar Documento
% Sistema: Hotel Perros

%-------------------------------------- COMIENZA descripción del caso de uso.

\begin{UseCase}{CU13}{Eliminar Documento}{
    Permite al Propietario eliminar un documento digital asociado a su mascota cuando este ya no es vigente, fue subido por error o se desea liberar espacio.
}
    \UCitem{Versión}{\color{Gray}1.0}
    \UCitem{Autor}{\color{Gray}Axel}
    \UCitem{Supervisa}{\color{Gray}[Nombre del Supervisor]}
    \UCitem{Actor}{\hyperlink{Propietario}{Propietario}}
    \UCitem{Propósito}{Mantener actualizado y organizado el repositorio digital de la mascota, eliminando archivos obsoletos o erróneos.}
    \UCitem{Entradas}{
        \begin{itemize}
            \item Selección del documento a eliminar (desde la lista).
            \item Confirmación de eliminación.
        \end{itemize}
    }
    \UCitem{Origen}{Mouse / Pantalla táctil}
    \UCitem{Salidas}{Mensaje de confirmación de eliminación exitosa, Lista de documentos actualizada.}
    \UCitem{Destino}{Pantalla}
    \UCitem{Precondiciones}{
        \begin{itemize}
            \item El propietario ha iniciado sesión.
            \item El propietario se encuentra visualizando la lista de documentos de una mascota.
            \item Existe al menos un documento registrado que pertenece a la mascota del propietario.
        \end{itemize}
    }
    \UCitem{Postcondiciones}{
        \begin{itemize}
            \item El registro del documento se elimina de la base de datos.
            \item El archivo físico se elimina del sistema de almacenamiento del servidor.
        \end{itemize}
    }
    \UCitem{Errores}{
        \begin{itemize}
            \item \textbf{E1}: El documento no existe o ya fue eliminado (Error 404).
            \item \textbf{E2}: El usuario no tiene permiso para eliminar este documento (Error 403).
            \item \textbf{E3}: Error interno al intentar borrar el archivo físico.
        \end{itemize}
    }
    \UCitem{Tipo}{Secundario}
\end{UseCase}

%--------------------------------------
% Trayectoria Principal
%--------------------------------------
\begin{UCtrayectoria}
    \UCpaso[\UCactor] Identifica el documento que desea borrar en la lista y oprime el botón \IUbutton{Eliminar} (icono de bote de basura o tache).
    \UCpaso Muestra el mensaje de confirmación {\bf MSG1-}``¿Está seguro que desea eliminar este documento permanentemente?''.
    \UCpaso[\UCactor] Oprime el botón \IUbutton{Sí, eliminar}.
    \UCpaso Solicita al servidor la eliminación del documento mediante su identificador único.
    \UCpaso Verifica que el documento exista y pertenezca al propietario autenticado.
    \UCpaso Elimina el archivo físico del servidor y su registro en la base de datos.
    \UCpaso Retorna una confirmación de éxito al cliente.
    \UCpaso Muestra el mensaje {\bf MSG2-}``Documento eliminado exitosamente''.
    \UCpaso Actualiza la lista de documentos, retirando el elemento eliminado.
\end{UCtrayectoria}

%--------------------------------------
% Trayectorias Alternativas
%--------------------------------------

% Trayectoria A: Cancelar
\begin{UCtrayectoriaA}{A}{El propietario cancela la eliminación}
    \UCpaso[\UCactor] Oprime el botón \IUbutton{Cancelar} en el mensaje de confirmación.
    \UCpaso Cierra el cuadro de diálogo sin realizar ninguna acción en el servidor.
    \UCpaso Continúa mostrando la lista de documentos original.
\end{UCtrayectoriaA}

% Trayectoria B: Documento no encontrado (Concurrencia)
\begin{UCtrayectoriaA}{B}{El documento ya no existe (Error 404)}
    \UCpaso El servidor intenta eliminar el documento pero no encuentra el registro (quizás fue eliminado desde otra sesión).
    \UCpaso Muestra el mensaje de error {\bf MSG-Error} ``Documento no encontrado''.
    \UCpaso Actualiza la lista automáticamente para reflejar el estado real (el documento desaparece).
\end{UCtrayectoriaA}

% Trayectoria C: Error de Permisos
\begin{UCtrayectoriaA}{C}{Intento de acceso no autorizado (Error 403)}
    \UCpaso El servidor detecta que el documento no pertenece al usuario actual.
    \UCpaso Muestra el mensaje de error {\bf MSG-Error} ``No autorizado para realizar esta acción''.
    \UCpaso Registra el intento de violación de seguridad.
\end{UCtrayectoriaA}